Uns electrons que es mouen horitzontalment travessen un selector de velocitats format per un camp magnètic de 0,040~T dirigit cap avall i un camp elèctric de 250~V/m perpendicular al camp magnètic i a la direcció de moviment dels electrons.

\figura[0.55]{fig1.png}

\begin{apartat}{1}
Dibuixeu i anomeneu les forces que actuen damunt l'electró quan és dins del selector de velocitats. Calculeu la velocitat dels electrons que travessaran el selector sense desviar-se.
\end{apartat}

\begin{apartat}{1}
Dins del selector un electró té una velocitat $\vec{v} = 1,25\times 10^{4}\,\vec{\imath}\ \mathrm{m\,s^{-1}}$ en el moment en què es desactiva el camp elèctric sense modificar el camp magnètic. Indiqueu la freqüència de rotació, el radi, el pla de gir i el sentit de gir del moviment circular uniforme d'aquest electró.
\end{apartat}

\dades{%
  $Q_\text{electró} = -1,60\times 10^{-19}\ \mathrm{C}$\\
  $m_\text{electró} = 9,11\times 10^{-31}\ \mathrm{kg}$}
\nota{Considereu negligible l'efecte de la força gravitatòria.}
